\phantomsection
\section*{Глава 1. Анализ предметной области}
\addcontentsline{toc}{section}{Глава 1. Анализ предметной области}
\stepcounter{section}

\subsection{Обзор существующих систем управления учебным процессом}

На современном рынке программного обеспечения для управления учебным процессом
сложились две основные категории решений. Универсальные системы управления
обучением (Learning Management Systems, LMS) - Moodle, iSpring Learn, Google
Classroom - ориентированы на дистанционное обучение и управление учебным контентом,
однако не предназначены для управления расписанием очных занятий и записью
на них. CRM-системы для образовательных организаций - Bitrix24, AmoCRM -
решают задачи учёта клиентов и продаж абонементов, но не предоставляют
функциональности для отслеживания успеваемости и работы с учебными материалами.

\textbf{Moodle} - наиболее распространённая открытая LMS, широко применяемая
в университетах и школах. Обеспечивает богатый набор инструментов для создания
курсов, тестирования и общения. Поддерживает ролевое разграничение доступа.
Среди недостатков с точки зрения языкового центра - отсутствие встроенного
модуля записи на конкретные занятия с ограничением числа мест, громоздкий
интерфейс, требующий значительных ресурсов на настройку и сопровождение.
Расписание реализуется только через сторонние плагины.

\textbf{Google Classroom} - бесплатный облачный сервис для организации учебного
процесса. Прост в освоении, тесно интегрирован с экосистемой Google (Диск,
Формы, Meet). Отсутствует механизм записи студентов на занятия с ограничением
вместимости; расписание не предусмотрено на уровне платформы. Ролевая модель
ограничена двумя ролями (учитель, ученик) и не поддерживает многоуровневое
администрирование.

\textbf{iSpring Learn} - российская коммерческая LMS, ориентированная на
корпоративное обучение. Предоставляет развитые инструменты аналитики и создания
интерактивных курсов. Регистрация на занятия и управление расписанием выходят
за рамки основного функционала; платформа не предназначена для небольших
образовательных организаций в силу высокой стоимости лицензии.

\textbf{Bitrix24} - многофункциональная CRM-система с образовательным модулем.
Позволяет вести базу клиентов, управлять продажами и отслеживать коммуникации.
Возможности отслеживания успеваемости студентов и работы с учебными материалами
минимальны; система ориентирована на бизнес-процессы, а не на организацию
учебных занятий.

Сравнительный анализ рассматриваемых платформ и разрабатываемой системы приведён
в таблице~\ref{tab:competitors}.

\begin{table}[ht!]
    \centering
    \caption{Сравнительный анализ систем управления учебным процессом}
    \label{tab:competitors}
    \small
    \begin{tabularx}{\textwidth}{|p{2.6cm}|X|p{1.7cm}|p{1.9cm}|p{2.0cm}|p{1.8cm}|}
        \hline
        \textbf{Платформа} &
        \textbf{Назначение} &
        \textbf{Запись на занятия} &
        \textbf{Расписание} &
        \textbf{Успевае- мость} &
        \textbf{Развёрты- вание} \\ \hline
        Moodle &
        Универсальная LMS для дистанционного обучения &
        Только через плагины &
        Только через плагины &
        Тесты и задания &
        Self-hosted / облако \\ \hline
        Google Classroom &
        Облачный сервис организации занятий &
        Отсутству- ет &
        Отсутству- ет &
        Оценки за задания &
        Только облако (Google) \\ \hline
        iSpring Learn &
        Корпоративная LMS &
        Отсутству- ет &
        Ограничено &
        Аналитика курсов &
        SaaS (платная) \\ \hline
        Bitrix24 &
        CRM с образовательным модулем &
        Через CRM-воронку &
        Календарь событий &
        Отсутствует &
        SaaS / self-hosted \\ \hline
        \textbf{Разрабатывае- мая система} &
        Веб-приложение языкового центра &
        Встроена, с лимитом мест &
        Встроено, с фильтрами &
        Оценки и посещаемость &
        Docker, self-hosted \\ \hline
    \end{tabularx}
\end{table}

Таким образом, ни одно из рассмотренных решений не обеспечивает одновременно:
запись студентов на конкретные занятия с контролем вместимости, управление
расписанием, отслеживание успеваемости и распределение учебных материалов в
рамках единого приложения, адаптированного для небольшого языкового центра.
Это подтверждает целесообразность разработки специализированного решения.

\subsection{Анализ технологий веб-разработки}

Выбор технологического стека определялся требованиями к реализации ролевой модели,
управлению расписанием и записями, агрегации данных об успеваемости и
контейнеризации приложения. Ниже приведено последовательное сравнение альтернатив
по каждому уровню архитектуры.

\textbf{Выбор серверного фреймворка.} Для реализации серверной части
рассматривались четыре фреймворка: Flask, Django, FastAPI и Express.js. Их
сравнение приведено в таблице~\ref{tab:frameworks}.

\begin{table}[ht!]
    \centering
    \caption{Сравнение серверных фреймворков}
    \label{tab:frameworks}
    \small
    \begin{tabularx}{\textwidth}{|p{2.3cm}|p{1.8cm}|X|X|p{2.2cm}|}
        \hline
        \textbf{Фреймворк} & \textbf{Язык} & \textbf{Преимущества} & \textbf{Недостатки} & \textbf{Решение} \\ \hline
        Flask &
        Python &
        Микрофреймворк с минимальными накладными расходами; гибкий выбор компонентов;
        богатая экосистема расширений (Flask-Login, Flask-WTF, Flask-Migrate, SQLAlchemy) &
        Нет встроенных ORM и панели администратора; всё настраивается вручную &
        \textbf{Выбран} \\ \hline
        Django &
        Python &
        «Батарейки в комплекте»: ORM, admin-панель, аутентификация, формы, миграции &
        Избыточная архитектура для проекта среднего масштаба; жёсткая структура &
        Не выбран \\ \hline
        FastAPI &
        Python &
        Асинхронная обработка; автоматическая OpenAPI-документация; типизация через Pydantic &
        Требует отдельного фронтенд-приложения; меньше учебных материалов &
        Не выбран \\ \hline
        Express.js &
        JavaScript &
        Высокая производительность; единый язык с фронтендом; широкое сообщество &
        JavaScript-экосистема вместо Python; менее удобная работа с формами &
        Не выбран \\ \hline
    \end{tabularx}
\end{table}

Flask выбран как микрофреймворк, позволяющий гибко формировать архитектуру
под конкретные требования. Расширения
Flask-Login, Flask-WTF и Flask-Migrate закрывают потребности в аутентификации,
защите форм и управлении миграциями без избыточности Django.

\textbf{Выбор системы управления базами данных.} Для хранения данных
рассматривались реляционные и нереляционные СУБД. Их сравнение приведено
в таблице~\ref{tab:databases}.

\begin{table}[ht!]
    \centering
    \caption{Сравнение систем управления базами данных}
    \label{tab:databases}
    \small
    \begin{tabularx}{\textwidth}{|p{2.5cm}|p{2.8cm}|X|X|p{2.2cm}|}
        \hline
        \textbf{СУБД} & \textbf{Тип} & \textbf{Преимущества} & \textbf{Недостатки} & \textbf{Решение} \\ \hline
        PostgreSQL &
        Реляционная &
        ACID-транзакции; сложные запросы и агрегации; поддержка составных ограничений
        уникальности; открытый код &
        Требует настройки отдельного сервера &
        \textbf{Выбрана} \\ \hline
        MySQL &
        Реляционная &
        Широкая распространённость; высокая производительность на операциях чтения &
        Исторически слабая поддержка транзакций; меньше аналитических функций &
        Не выбрана \\ \hline
        SQLite &
        Реляционная &
        Не требует сервера; нулевая настройка; удобна для разработки &
        Ограниченная конкурентная запись; не пригодна для продакшена при нагрузке &
        Не выбрана \\ \hline
        MongoDB &
        Документная (NoSQL) &
        Гибкая схема без миграций; горизонтальное масштабирование &
        Слабая поддержка JOIN; более сложные агрегации для расчёта успеваемости &
        Не выбрана \\ \hline
    \end{tabularx}
\end{table}

PostgreSQL выбрана как реляционная СУБД с полной поддержкой транзакций,
составных уникальных ограничений (например, одна запись студента на занятие)
и агрегатных функций, необходимых для расчёта среднего балла и статистики
посещаемости. Работа с базой данных осуществляется через
ORM-библиотеку SQLAlchemy, обеспечивающую декларативное
описание моделей и систему миграций через Flask-Migrate (Alembic).

\textbf{Выбор подхода к формированию фронтенда.} Для рендеринга пользовательских
страниц рассматривались три подхода: серверный рендеринг (SSR), одностраничные
приложения (SPA) и гибридный HTML-first. Их сравнение приведено
в таблице~\ref{tab:frontend}.

\begin{table}[ht!]
    \centering
    \caption{Сравнение подходов к разработке фронтенда}
    \label{tab:frontend}
    \small
    \begin{tabularx}{\textwidth}{|p{2.5cm}|p{2.3cm}|X|X|p{2.2cm}|}
        \hline
        \textbf{Подход} & \textbf{Технология} & \textbf{Преимущества} & \textbf{Недостатки} & \textbf{Решение} \\ \hline
        SSR-шаблоны &
        Jinja2 (Flask) &
        Встроен во Flask; не требует отдельного frontend-проекта и сборки;
        прямая интеграция с моделями данных; низкий порог входа &
        Ограниченная клиентская интерактивность без дополнительного JavaScript &
        \textbf{Выбран} \\ \hline
        SPA &
        React / Vue.js &
        Высокая интерактивность; компонентный подход; богатый UX &
        Требует отдельного REST API, разделения репозитория и дополнительного
        процесса сборки и деплоя &
        Не выбран \\ \hline
        HTML-first &
        HTMX &
        Минимум JavaScript; частичная замена фрагментов страницы без перезагрузки &
        Ограниченный набор UI-паттернов; небольшое сообщество &
        Не выбран \\ \hline
        Гибридный SSR &
        Next.js &
        SEO; производительность; поддержка React-компонентов &
        JavaScript/Node.js стек; избыточная сложность для данного проекта &
        Не выбран \\ \hline
    \end{tabularx}
\end{table}

Jinja2 выбран как шаблонизатор, встроенный во Flask,
обеспечивающий серверный рендеринг HTML-страниц без подключения отдельного
frontend-фреймворка. Такой подход соответствует масштабу проекта и устраняет
необходимость раздельного деплоя backend-API и frontend-приложения.

\textbf{Выбор инструмента контейнеризации.} Для обеспечения воспроизводимости
среды развёртывания рассматривались Docker и виртуальные машины (VM).
Docker выбран в качестве платформы контейнеризации: контейнеры
обеспечивают изоляцию приложения и базы данных, а Docker Compose позволяет
описать многосервисную конфигурацию (web + db) в одном файле и запустить
всё окружение одной командой. В отличие от виртуальных машин, контейнеры
не требуют полноценной гостевой ОС и потребляют значительно меньше ресурсов.

Ключевые технологические решения итоговой архитектуры:

\begin{itemize}
    \item \textbf{аутентификация:} сессионная модель на основе Flask-Login с
          хешированием паролей через \texttt{werkzeug.security}; защита форм
          через Flask-WTF (CSRF-токены);
    \item \textbf{миграции БД:} Flask-Migrate (Alembic) обеспечивает
          версионирование схемы базы данных и применение изменений без потери данных;
    \item \textbf{загрузка файлов:} учебные материалы сохраняются в файловой
          системе контейнера; имена файлов генерируются через UUID во избежание
          коллизий; выгрузка реализована через \texttt{flask.send\_file}.
\end{itemize}

\subsection{Формулировка цели и задач работы}

На основе проведённого обзора существующих систем (см. подраздел 1.1) и
анализа технологического стека (см. подраздел 1.2) сформулирована цель работы.

\textbf{Цель работы} - разработка веб-приложения для центра изучения иностранных
языков, обеспечивающего запись студентов на занятия, управление расписанием
и отслеживание успеваемости в рамках единой системы с разграничением доступа
по ролям.

Для достижения цели необходимо решить следующие задачи:
\begin{enumerate}
    \item Провести анализ существующих систем управления учебным процессом и
          обосновать выбор технологического стека.
    \item Спроектировать модель данных для сущностей системы: пользователь,
          курс, занятие, запись, оценка, учебный материал, уведомление.
    \item Реализовать систему аутентификации и ролевого разграничения доступа
          (администратор, преподаватель, студент).
    \item Реализовать систему записи студентов на занятия с контролем вместимости
          и управление недельным расписанием с фильтрацией.
    \item Реализовать CRUD-интерфейсы для управления курсами, занятиями
          и пользователями.
    \item Реализовать загрузку и выгрузку учебных материалов с хранением
          файлов на сервере.
    \item Реализовать систему отслеживания успеваемости: выставление оценок
          и посещаемости преподавателем, агрегирование показателей по студенту.
    \item Реализовать систему внутренних уведомлений с отображением в реальном
          времени и экспортом оценок в формате CSV.
    \item Выполнить контейнеризацию приложения с помощью Docker и оформить
          отчётную документацию.
\end{enumerate}
